This Online Appendix provides additional information on the staggered rollout of the \emph{Plan Cuadrante} program.


\emph{Plan Cuadrante} was launched in 1998, initially as a pilot program in a few municipalities in the southern part of the city of Santiago. This initial phase aimed to test the plan's implementability. The six-month pilot revealed promising outcomes, which led to its broader implementation despite the need for additional resources and personnel.

By 2000, the \emph{Plan Cuadrante} had expanded to 44 municipalities in Santiago, covering approximately 96.1\% of the city’s population. This expansion was based on expert recommendations and the division of urban areas into manageable patrol zones. The plan was then implemented gradually throughout urban municipalities in Chile. The staggered expansion of the \emph{Plan Cuadrante} beyond Santiago was, in principle, meant to follow a prioritization based on a weighted index of victimization, police resource deficits, drug prevalence, unemployment, and demand for police services, with annual entry determined by available budget. The weights assigned to each component changed on a yearly basis and, in practice, Budget Office evaluations show that adoption diverged from this ranking for undocumented reasons, so the rule operated more as a loose guideline than as a binding assignment mechanism \citep{DIPRES2014}.\footnote{We requested information from Carabineros about the index, its inputs and weights but we were not granted systematic access to this information.}


The plan saw significant growth between 2002 and 2006. In 2002, it extended beyond Santiago into the V and VIII Regions, adding ten new municipalities. The expansion continued from 2003 to 2006, with 18 additional municipalities incorporated into the plan. By 2006, the \emph{Plan Cuadrante} covered a total of 72 municipalities across Chile. Between 2007 and 2009, the plan expanded to include 28 more municipalities, reaching the 100 municipalities covered by the \emph{Plan Cuadrante}. The \emph{Plan Cuadrante} faced significant challenges in 2010 due to the earthquake that struck Chile. The disaster necessitated a reallocation of resources to support recovery efforts, temporarily halting the plan’s expansion. Despite this setback, the \emph{Plan Cuadrante} resumed its implementation in 2011. Between 2011 and 2013, the plan was expanded to include an additional 50 municipalities, bringing the total to 150 across Chile. Tables \ref{tab:sum_statsA2} and \ref{tab:sum_statsA3} present summary statistics at the household and at the municipality level of the municipalities implementing \emph{Plan Cuadrante} in the period that we study. 

\vspace{5mm}
\begin{table}[H]
\centering
\caption{Implementation of \emph{Plan Cuadrante} by Year}
\label{tab:implementation}
\begin{threeparttable}
\begin{tabular}{l c >{\raggedright\arraybackslash}p{10.5cm}}
\toprule
Year & N & Municipalities \\
\midrule
1998--2000 & 44  & Metropolitan Region of Santiago (44 municipalities) \\[3pt]
2002 & 10  & Chiguayante, Concepción, Hualpén, Quilpué, San Antonio, San Pedro de la Paz, Talcahuano, Valparaíso, Villa Alemana, Viña del Mar \\[3pt]
2003 & 2   & Padre Las Casas, Temuco \\[3pt]
2004 & 2   & Antofagasta, Copiapó \\[3pt]
2005 & 7   & Alto Hospicio, Arica, Calama, Curicó, Iquique, Linares, Talca \\[3pt]
2006 & 7   & Coquimbo, La Serena, Los Andes, Ovalle, Rancagua, San Felipe, San Fernando \\[3pt]
2007 & 7   & Chillán, Chillán Viejo, Los Ángeles, Osorno, Puerto Montt, Punta Arenas, Valdivia \\[3pt]
2008 & 9   & Ancud, Castro, Constitución, Coronel, Lota, Peñaflor, Quillota, Rengo, Villarrica \\[3pt]
2009 & 12  & Angol, Aysén, Calera, Concón, Coyhaique, La Unión, Padre Hurtado, Parral, Penco, San Carlos, Tomé, Vallenar \\[3pt]
2011 & 17  & Arauco, Cañete, Cartagena, Cauquenes, Curanilahue, La Ligua, Limache, Molina, Nueva Imperial, Pucón, Puerto Varas, Río Bueno, San Clemente, San Javier, San Vicente, Santa Cruz, Victoria \\[3pt]
2012 & 20  & Cabrero, Calbuco, Casablanca, Chimbarongo, Collipulli, Curacaví, El Monte, Graneros, Illapel, Isla de Maipo, Lautaro, Machalí, Mostazal, Mulchén, Nacimiento, Natales, Panguipulli, Quellón, Quintero, Tocopilla \\[3pt]
2013 & 13  & Bulnes, Laja, Las Cabras, Lebu, Llaillay, Loncoche, Longaví, Los Álamos, Maule, Requínoa, Teno, Vicuña, Yumbel \\
\midrule
Total & 150 & \\
\bottomrule
\end{tabular}
\begin{tablenotes}[flushleft]
\footnotesize
\item Note: \textit{N} is the number of municipalities that adopted \emph{Plan Cuadrante} in each year. Data shared by Carabineros de Chile.
\end{tablenotes}
\end{threeparttable}
\end{table}


\begin{table}[H]
\centering
\caption{Summary Statistics -- ENUSC sample}
\label{tab:sum_statsA2}
\begin{threeparttable}
\begin{tabular}{l r r r r r}
\toprule
 & Mean & SD & Min & Max & N \\
\midrule
\multicolumn{6}{@{}l}{\textit{Household}} \\
\quad Socioeconomic status               & 3.58  & 0.73  & 1    & 5    & 99,298 \\
\quad Victimization                    & 0.26  & 0.44  & 0    & 1    & 99,298 \\
\quad Crime perception (neighborhood)  & 0.38  & 0.49  & 0    & 1    & 95,598 \\
\quad Crime perception (municipality)  & 0.66  & 0.48  & 0    & 1    & 96,794 \\
\quad Crime perception (country)       & 0.79  & 0.41  & 0    & 1    & 98,525 \\
\quad Trust in the police              & 0.44  & 0.50  & 0    & 1    & 51,445 \\
\quad Investment in security measures  & 0.21  & 0.40  & 0    & 1    & 99,298 \\
\quad \emph{Plan Cuadrante} (year)                  & 2007  & 1.85  & 2005 & 2012 & 97,847 \\
\addlinespace
\multicolumn{6}{@{}l}{\textit{Household member interviewed}} \\
\quad Female                      & 0.55  & 0.50  & 0  & 1  & 99,298 \\
\quad Age                         & 43.87 & 17.95 & 15 & 99 & 99,298 \\
\quad Primary school attendance   & 0.98  & 0.14  & 0  & 1  & 99,298 \\
\quad Secondary school attendance & 0.72  & 0.45  & 0  & 1  & 99,298 \\
\quad University attendance       & 0.20  & 0.40  & 0  & 1  & 99,298 \\
\bottomrule
\end{tabular}
\begin{tablenotes}[flushleft]
\footnotesize
\item Note: This table presents summary statistics from the ENUSC sample, covering the 47 municipalities used in the survey-based analysis (46 that adopted \emph{Plan Cuadrante} between 2005 and 2012 plus one never-treated municipality). Socioeconomic status is a 1 to 5 measure produced by ENUSC surveys and based on self-reported socioeconomic strata, being 1 the highest ranked and 5 the lowest. Crime perception is equal to 1 if the interviewed person reports that crime has increased in the last 12 months. The \emph{Plan Cuadrante} (year) row is computed over treated municipalities only.
\end{tablenotes}
\end{threeparttable}
\end{table}

\begin{table}[H]
\centering
\caption{Summary Statistics -- Administrative municipality level sample}
\label{tab:sum_statsA3}
\begin{threeparttable}
\begin{tabular}{l r r r r r}
\toprule
 & Mean & SD & Min & Max & N \\
\midrule
\multicolumn{6}{@{}l}{\textit{Panel A. Reported-crime sample (292 municipalities)}} \\
\quad Reported crime (per 100k pop.)            & 1,732.39 & 1,036.68 & 0.00   & 11,434.40 & 4,371 \\
\quad Arrests (per 100k pop.)                   & 332.07   & 271.75   & 0.00   & 1,712.30  & 3,504 \\
\quad Poverty rate (\%)                         & 25.74    & 9.72     & 5.65   & 58.33     & 3,720 \\
\quad Pop.\ density (pop./km\textsuperscript{2}) & 63.51   & 140.47   & 0.09   & 1,109.58  & 4,065 \\
\quad \emph{Plan Cuadrante} (year)                     & 2010     & 2.82     & 2003   & 2013      & 1,440 \\
\addlinespace
\multicolumn{6}{@{}l}{\textit{Panel B. Arrests sample (281 municipalities)}} \\
\quad Reported crime (per 100k pop.)            & 1,684.74 & 1,014.67 & 0.00   & 11,434.40 & 4,209 \\
\quad Arrests (per 100k pop.)                   & 313.21   & 251.23   & 0.00   & 1,712.30  & 3,372 \\
\quad Poverty rate (\%)                         & 26.01    & 9.72     & 5.65   & 58.33     & 3,570 \\
\quad Pop.\ density (pop./km\textsuperscript{2}) & 57.90   & 127.30   & 0.09   & 1,109.58  & 3,900 \\
\quad \emph{Plan Cuadrante} (year)                     & 2010     & 2.26     & 2006   & 2013      & 1,275 \\
\bottomrule
\end{tabular}
\begin{tablenotes}[flushleft]
\footnotesize
\item Note: The table reports municipality-level summary statistics for the two administrative estimation samples used in the paper. Panel A covers the 292 municipalities in the reported-crime sample (municipalities adopting \emph{Plan Cuadrante} from 2003 onward plus never-treated); Panel B the 281 municipalities in the arrests sample (adopters from 2006 onward plus never-treated --- municipalities adopting in 2005 are excluded because arrest data begin that year, leaving no pre-treatment period). Reported crime and arrests are measured per 100,000 inhabitants over 2002--2016 and 2005--2016, respectively. Minimum reported crime and arrests of zero correspond to tiny, remote never-treated comunas (e.g., Juan Fernández, Cabo de Hornos). The \emph{Plan Cuadrante} (year) row is computed over treated municipalities only. Data from the Undersecretary for Crime Prevention.
\end{tablenotes}
\end{threeparttable}
\end{table}

\clearpage
